\SourceLeaf{347}

\begin{center}
[II.]

\textbf{LEHRSÄTZE UND AUFGABEN}

\textbf{DIE COMPLEXEN ZAHLEN BETREFFEND.}
\end{center}

\Needspace{7\baselineskip}1. $\dfrac{a}{b-\rho}$ stellt, wenn $a$, $b$ bestimmte Zahlen, $\rho$ eine
unbestimmte cyklische Zahl bedeutet, einen Kreis vor,
\[
\dfrac{ab}{bb-\beta\beta}
\]
den Mittelpunkt, wenn $\beta$ der in $b$ enthaltene cyklische Factor,
\[
\dfrac{a}{bb-\beta\beta}
\]
(oder der Modulus dieser Grösse) den Halbmesser des Kreises.

\Needspace{7\baselineskip}2. (vor 1.) Ist wie oben $\beta$ der in $b$ enthaltene cyklische Factor, so
haben
\[
b + \beta\rho \quad \text{und} \quad b + \dfrac{\beta}{\rho}
\]
einerlei Modulus, imgleichen
\[
b + \beta\rho \quad \text{und} \quad \dfrac{\rho\beta\beta}{b\rho+\beta}
\]
einerlei Clise.

\Needspace{7\baselineskip}3. $\rho - \beta$ und $i\sqrt{\rho\beta}$ haben gleiche Clise.

\Needspace{7\baselineskip}4. $\Omega = \dfrac{a\rho+b}{1-\rho}$ stellt eine gerade Linie vor. Man hat
nemlich
\[
2\Omega = \dfrac{(a+b)(1+\rho) - (a-b)(1-\rho)}{1-\rho}
\]
\[
= (a+b)t - (a-b),
\]
wenn
\[
t = \dfrac{1+\rho}{1-\rho}
\]
gesetzt wird, wodurch also $t$ eine unbestimmte imaginäre Zahl wird. Es ist
ferner
\[
\dfrac{2\Omega}{a+b} = t - \dfrac{a-b}{a+b}.
\]

Derjenige Werth dieser Grösse, dessen Modulus ein Minimum ist, ist offenbar
nichts anderes als der reelle Theil von $\dfrac{a-b}{a+b}$.

Ist demnach $k$ der reelle Theil von $\dfrac{a-b}{a+b}$, so ist der
gesuchte Werth von $\Omega$
\[
= \tfrac{1}{2} k (a+b).
\]
